% !TeX root = ../../Book2.tex
\chapter{Schwartz 空间与温和分布}
\label{b2:ch:05}
\BookChapterNumber{b2:ch:05}
% 资料：Hunter2001Applied，作者公开ch11.pdf，实际渲染核读291、294、303--304页。
% Schwartz半范数、双线性配对、归一化转换与正定Gauss公式分别核对。
% 本章自足证明Hermite完备性、S'自同构和各维负Laplace基本解。
\begin{chapterabstract}
Fourier 变换要容纳常函数、点质量和可积函数以外的对象，需要一个在变换下保持不变的测试空间。本章先构造 Schwartz 空间，证明半范数完备性及微分、多项式乘法、卷积和变换的连续性，并在选读专题中以 Hermite 函数展示变换的离散特征结构。随后把变换转移到连续对偶，得到温和分布上的统一理论。最后将常系数微分方程化为频率乘法，构造带正零阶项及各维负 Laplace 算子的基本解。
\end{chapterabstract}
\begin{learninggoals}
\begin{itemize}
\item\label{b2:item:05-goal-schwartz}使用具体半范数检查快速下降、收敛和连续操作，区分 Schwartz 收敛与共同紧支集的测试函数收敛。
\item\label{b2:item:05-goal-hermite}选读进阶目标：在\subsecref{b2:subsec:050204}中由 Gauss 函数构造 Hermite 系，理解其正交性、完备性和 Fourier 特征值。
\item\label{b2:item:05-goal-tempered}用连续泛函定义温和分布，保持双线性配对与既有微分符号一致。
\item\label{b2:item:05-goal-examples}计算 Dirac 测度、多项式、导数与齐次分布的变换，判断普通积分何时不能再用。
\item\label{b2:item:05-goal-fundamental}将常系数算子转为多项式符号，计算不同维数的基本解，并说明原点奇性和边界定向在证明中的作用。
\end{itemize}
\end{learninggoals}

一单位质量集中在原点时，普通函数无法同时表示它的总质量和无限集中程度。另一个看似相反的对象是常函数 \(1\)：它均匀铺满空间，却应只含零频率。两者的 Fourier 变换应彼此对应；要把这句话变成计算规则，需要先找到对二者都可测试、又在变换下保持不变的函数空间。本章构造这个空间，然后用点质量的方程 \(-\Delta E=\delta_0\) 说明广义对象怎样产生可用的解。

\ProseParagraphBreak

本章使用\chapref{b2:ch:03}的变换公式与反演、\chapref{b2:ch:04}的平方能量理论，以及\chapref{b2:ch:01}的分部积分和 Green 公式。第一卷定义 20.2.1 的普通分布作为比较对象：其测试函数紧支集，故不要求控制无穷远；温和分布则须连续作用于全部 Schwartz 函数。二者的范围不同，不能只因名称相似就互相替换。

\ProseParagraphBreak

\secref{b2:sec:0501}和\secref{b2:sec:0502}讨论 Schwartz 函数及其变换；\secref{b2:sec:0503}通过对偶配对定义温和分布上的 Fourier 变换；\secref{b2:sec:0504}用它研究常系数方程。John K. Hunter 与 Bruno Nachtergaele 的《Applied Analysis》\cite[Chapter 11]{Hunter2001Applied}介绍测试函数与温和分布。Lars Hörmander 的《The Analysis of Linear Partial Differential Operators I: Distribution Theory and Fourier Analysis》\cite{Hrmander2003Analysis}详细论述一般分布与常系数方程理论，可在学完本章后按具体问题查阅。

% !TeX root = ../../Book2.tex
\section{一阶线性方程}
\label{b2:sec:0701}

给定某一时刻的浓度分布，若介质只作平移而没有扩散，后来的图形应当保持原来的形状。我们先用方程描述这种平移，再沿直线求出解，由此介绍特征线法。本章把空间变量写在前面，时间变量写在后面，即 \(u=u(x,t)\)。向量场表示速度，函数值既可能表示随粒子携带的量，也可能表示单位体积中的密度；这两种含义将在\secref{b2:sec:0703}区分。

\subsection{输运方程}
\label{b2:subsec:070101}

设 \(b\in\mathbb R^d\) 为固定速度。经过初始位置 \(a\) 的粒子在时刻 \(t\) 位于 \(a+tb\)。若沿这一运动携带的数值不变，则链式法则给出
\[
 \frac{\mathrm d}{\mathrm dt}u(a+tb,t)
 =u_t(a+tb,t)+b\cdot\nabla_xu(a+tb,t)=0.
\]
这产生\term[shuyun-fangcheng]{输运方程}[transport equation]
\begin{equation}\label{b2:eq:07-constant-transport}
 u_t+b\cdot\nabla_xu=0.
\end{equation}
其中 \(u_t\) 衡量固定观察点的变化，\(b\cdot\nabla_xu\) 衡量移动造成的变化。二者可以分别非零而恰好相消。例如一维函数 \(u(x,t)=\sin(x-bt)\) 的每个波峰向右或向左平移，方向取决于 \(b\) 的符号。

\ProseParagraphBreak

这里讨论的经典解属于 \(C^1\)，因此上述链式法则有逐点意义。本卷\chapref{b2:ch:04}的分布语言允许进一步处理不光滑对象，但本章先把几何和唯一性建立在光滑范围内；仅在明确写出测试函数恒等式时才转入弱解。

\subsection{常系数情形}
\label{b2:subsec:070102}

\begin{proposition}\label{b2:prop:07-translation-solution}
设 \(I\) 为包含 \(0\) 的区间，\(u\in C^1(\mathbb R^d\times I)\)。函数 \(u\) 满足\eqref{b2:eq:07-constant-transport}，当且仅当存在 \(g\in C^1(\mathbb R^d)\)，使
\[
 u(x,t)=g(x-tb).
\]
其中 \(g(x)=u(x,0)\)，因而初值唯一确定解。
\end{proposition}
\begin{proof}
固定 \((x,t)\)，取 \(a=x-tb\)。函数 \(s\mapsto u(a+sb,s)\) 的导数为零，从 \(0\) 积分到 \(t\) 得 \(u(x,t)=u(a,0)\)。反过来，对所给公式求导，有 \(u_t=-b\cdot\nabla g(x-tb)\)、\(\nabla_xu=\nabla g(x-tb)\)，故方程成立。时间为负时，积分的方向反转，论证不变。
\end{proof}

平移不会提高正则性：初值在哪里存在折角，后来就在相应平移点保留折角。对于任意 \(1\leq p\leq\infty\)，若初值还属于 \(L^p\)，平移不变性给出 \(\|u(\cdot,t)\|_p=\|g\|_p\)。这一恒等式没有衰减因子，与以后热方程的平滑效应不同。

\ProseParagraphBreak

不把某个变量预先叫作时间时，同样的方法适用于 \(\alpha u_x+\beta u_y=0\)，其中 \((\alpha,\beta)\ne(0,0)\)。取
\[
 r=\alpha x+\beta y,\quad s=\beta x-\alpha y,
\]
则变换的行列式为 \(-\alpha^2-\beta^2\ne0\)，算子化为 \((\alpha^2+\beta^2)\partial_r\)。所以在整个平面上，解恰为 \(F(\beta x-\alpha y)\)。若区域与一条直线的交集不连通，各个连通片上的常数未必由同一个全局函数 \(F\) 描述；局部坐标公式不能自动升级为任意区域上的全局表示。

\subsection{非齐次方程}
\label{b2:subsec:070103}

源项使粒子所携带的值随时间累积，零阶项则造成吸收或增长。先考虑 \(\lambda\in\mathbb R\) 为常数的情形。

\begin{theorem}\label{b2:thm:07-constant-source}
设 \(g\in C^1(\mathbb R^d)\)、\(f\in C^1(\mathbb R^d\times[0,T])\)。问题
\[
 u_t+b\cdot\nabla u+\lambda u=f,\quad u(x,0)=g(x)
\]
有唯一的 \(C^1\) 解，且
\begin{equation}\label{b2:eq:07-constant-source}
 u(x,t)=e^{-\lambda t}\cdot g(x-tb)
 +\int_0^t e^{-\lambda(t-s)}\cdot f\bigl(x-(t-s)b,s\bigr)\,\mathrm ds.
\end{equation}
闭时间区间上的光滑性均理解为能局部延伸到邻域，端点导数也可取相应单侧值。
\end{theorem}
\begin{proof}
沿 \(x=a+tb\) 令 \(v(t)=u(a+tb,t)\)，方程成为
\(v'+\lambda v=f(a+tb,t)\)。乘以 \(e^{\lambda t}\) 并从 \(0\) 积分，得到
\[
 v(t)=e^{-\lambda t}\cdot g(a)
 +\int_0^t e^{-\lambda(t-s)}\cdot f(a+sb,s)\,\mathrm ds.
\]
令 \(a=x-tb\) 即为表示公式，故任何经典解只能取此值。公式右侧在每个紧时空集合上的被积函数及一阶导数都有共同界，\secref{b2:sec:0103}的积分求导规则保证它属于 \(C^1\)。把它沿每条直线拉回，所得函数满足刚才的常微分方程；链式法则再给出原方程。取 \(t=0\) 得到初值。
\end{proof}

例如在一维取 \(b=1\)、\(\lambda=0\)、\(g=0\)、\(f(x,t)=x\)，有
\[
 u(x,t)=\int_0^t(x-t+s)\,\mathrm ds=tx-\frac{t^2}{2}.
\]
这里源项须在粒子过去所在的位置取值；若误把积分写成 \(\int_0^tf(x,s)\,\mathrm ds\)，会得到不满足方程的 \(tx\)。

\subsection{初值问题}
\label{b2:subsec:070104}

初值面 \(t=0\) 被每条轨道恰好穿过一次，因此可以给每条轨道一个独立数值。若改在与轨道相切的曲线上给值，存在性或唯一性就会失效。以 \(u_x+u_y=0\) 为例，沿直线 \(y=x\) 解必须恒定；指定 \(u(s,s)=s\) 没有解，而指定 \(u(s,s)=0\) 则允许所有满足 \(F(0)=0\) 的解 \(u(x,y)=F(x-y)\)。方程与数据之间的几何位置是初值问题的一部分。

\ProseParagraphBreak

常系数公式还直接给出数据稳定性。若 \(u,\widetilde u\) 对应 \(g,f\) 与 \(\widetilde g,\widetilde f\)，并且右侧范数有限，则
\[
 \|u(\cdot,t)-\widetilde u(\cdot,t)\|_\infty
 \leq e^{|\lambda|t}\left(\|g-\widetilde g\|_\infty
 +\int_0^t\|f(\cdot,s)-\widetilde f(\cdot,s)\|_\infty\,\mathrm ds\right).
\]
证明只需在\eqref{b2:eq:07-constant-source}中相减并取绝对值。下一节把直线改成由速度场决定的曲线；为使这种替换合法，先要证明曲线存在、彼此不发生碰撞，并且随初始点可微地变化。

% !TeX root = ../../Book2.tex
\section{特征曲线}
\label{b2:sec:0702}

变系数速度场已经给出粒子在每一点的瞬时运动方向，却没有直接给出整条轨道。先解常微分方程，再用轨道作坐标，是本节的构造顺序。以下 \(U\subseteq\mathbb R^d\) 为开集，\(b\in C^1(U;\mathbb R^d)\)；本节先处理不显含时间的速度场。

\subsection{特征常微分方程}
\label{b2:subsec:070201}

\begin{definition}\label{b2:def:07-characteristic}
若 \(X:I\rightarrow U\) 属于 \(C^1\)，并满足
\[
 X'(t)=b\bigl(X(t)\bigr),
\]
则称它为速度场 \(b\) 的\term[tezheng-quxian]{特征曲线}[characteristic curve]，也称特征线。其常微分方程称为\term[tezheng-changweifen-fangcheng]{特征常微分方程}[characteristic ordinary differential equation]。
\end{definition}

沿这样的曲线，复合函数的导数只涉及 \(u_t+b\cdot\nabla u\)。但在应用链式法则之前，要确认对给定初始位置确实有一条且只有一条曲线。下面给出局部构造，并保留稍后需要的定量控制。

\begin{theorem}[Picard--Lindelöf]\label{b2:thm:07-local-ode}
对每个 \(a_0\in U\)，存在 \(r>0\)、\(h>0\)，使每个 \(a\in B_r(a_0)\) 都有唯一的局部解
\[
 X'(t)=b\bigl(X(t)\bigr),\quad X(0)=a,
 \quad |t|\leq h.
\]
各解取值于同一个紧球。若 \(L\) 是该球上 \(\|Db\|\) 的上界，则
\begin{equation}\label{b2:eq:07-flow-lipschitz}
 |X(t,a)-X(t,\widetilde a)|\leq e^{L|t|}\cdot|a-\widetilde a|.
\end{equation}
局部唯一性指任意两条满足同一初值的解在共同定义区间上相等。
\end{theorem}
\begin{proof}
取闭球 \(\overline B_{2r}(a_0)\subseteq U\)，记其上 \(|b|\)、\(\|Db\|\) 的上界为 \(M,L\)。选 \(h>0\) 满足 \(Mh\leq r\)、\(Lh<1\)；上界为零的条件可略去。对固定 \(a\in B_r(a_0)\)，在连续曲线空间中考虑闭集
\[
 E_a=\{\,Y\in C([-h,h];\mathbb R^d)\mid
 \sup_{|t|\leq h}|Y(t)-a|\leq r\,\}.
\]
一致 Cauchy 曲线列逐点有极限，一致极限仍连续且满足闭球约束，所以 \(E_a\) 关于一致范数完备。定义
\[
 (TY)(t)=a+\int_0^t b\bigl(Y(s)\bigr)\,\mathrm ds.
\]
速度界保证 \(T(E_a)\subseteq E_a\)，中值估计保证
\(\|TY-TZ\|_\infty\leq Lh\cdot\|Y-Z\|_\infty\)。从常曲线 \(Y_0=a\) 迭代，连续两次之差不超过 \((Lh)^n\cdot\|Y_1-Y_0\|_\infty\)。几何级数收敛，故迭代一致收敛到 \(X\in E_a\)。算子连续性给出 \(TX=X\)；积分的被积函数连续，所以 \(X\) 可微并满足方程。两个不动点之间的距离不超过原距离的 \(Lh\) 倍，故相等。

为证明稳定估计，先记录一个积分不等式。若连续非负函数 \(q\) 满足 \(q(t)\leq A+L\int_0^tq(s)\,\mathrm ds\)，令 \(Q(t)=A+L\int_0^tq(s)\,\mathrm ds\)，则 \(Q'\leq LQ\)，从而 \((e^{-Lt}Q)'\leq0\)，得到 \(q(t)\leq Ae^{Lt}\)。此结论也覆盖 \(A=0\)。对两个解的积分方程作差即可用它得到\eqref{b2:eq:07-flow-lipschitz}；负时间用 \(s=-t\) 处理。任意两个局部解在某个更短区间上都落入同一个球，故先在那里相等。在共同区间内，相等的时刻构成闭集，局部唯一性又使它为开集，因而它充满包含初始时刻的整个区间。
\end{proof}

上述结果称为\term[picard-lindelof]{Picard--Lindelöf 定理}[Picard--Lindelöf theorem]，证明中的积分估计是\term[gronwall-budengshi]{Gronwall 不等式}[Gronwall inequality]的一个基本形式。只要求速度连续时，唯一性可能失败。例如 \(X'=2\sqrt{|X|}\)、\(X(0)=0\) 除零解外，对每个 \(a\geq0\) 还允许先在 \([0,a]\) 上停留，之后取 \(X(t)=(t-a)^2\)。本章采用 \(C^1\) 速度，正是为了排除这种轨道分叉。

\subsection{流映射}
\label{b2:subsec:070202}

把从 \(a\) 出发的解记为 \(\Phi_t(a)\)，称 \(a\mapsto\Phi_t(a)\) 为\term[liu-yingshe]{流映射}[flow map]。局部解借唯一性在重叠区间上拼接，得到经过 \(a\) 的最大解区间 \(I_a\)。记号 \(\Phi_t(a)\) 只在 \(t\in I_a\) 时使用；不能因为给了一个映射名称便假定它对所有时间存在。
\symidx{\(\Phi_t\)：自治速度场的流映射}

\begin{proposition}\label{b2:prop:07-flow-composition}
凡各项有定义，流满足
\[
 \Phi_0(a)=a,\quad
 \Phi_{t+s}(a)=\Phi_t\bigl(\Phi_s(a)\bigr),\quad
 \Phi_{-t}\bigl(\Phi_t(a)\bigr)=a.
\]
固定 \(t\) 时，\(\Phi_t\) 在其定义域上单射，其逆为相应的 \(\Phi_{-t}\)。
\end{proposition}
\begin{proof}
曲线 \(r\mapsto\Phi_{s+r}(a)\) 与 \(r\mapsto\Phi_r(\Phi_s(a))\) 在 \(r=0\) 具有相同初值，且都满足自治方程。唯一性给出复合关系。令 \(r=-s\) 得逆关系，进而得到单射性。
\end{proof}

\begin{theorem}\label{b2:thm:07-flow-derivative}
流 \((t,a)\mapsto\Phi_t(a)\) 在其开定义域上属于 \(C^1\)。对初始点的导数 \(A(t,a)=D_a\Phi_t(a)\) 满足
\begin{equation}\label{b2:eq:07-variational-ode}
 A'(t,a)=Db\bigl(\Phi_t(a)\bigr)A(t,a),\quad A(0,a)=I.
\end{equation}
因此固定时间的流是其定义域到像集的 \(C^1\) 微分同胚。
\end{theorem}
\begin{proof}
先在上一构造的紧球和短时间内工作。式\eqref{b2:eq:07-flow-lipschitz}与 \(|\partial_t\Phi_t|\leq M\) 保证流联合连续。矩阵积分方程
\[
 A(t)=I+\int_0^t Db\bigl(\Phi_s(a)\bigr)A(s)\,\mathrm ds
\]
可用同样的迭代构造，且 \(\|A(t)\|\leq e^{L|t|}\)。现取小向量 \(v\)，令 \(\delta_v(t)=\Phi_t(a+v)-\Phi_t(a)\)。中值积分给出
\[
 b\bigl(\Phi_t(a+v)\bigr)-b\bigl(\Phi_t(a)\bigr)
 =Db\bigl(\Phi_t(a)\bigr)\delta_v(t)+r_v(t).
\]
导数 \(Db\) 在紧球上一致连续，而 \(|\delta_v(t)|\leq e^{Lh}|v|\)，所以
\(\sup_t|r_v(t)|/|v|\to0\)。将 \(\delta_v-Av\) 的两个积分方程相减，再用刚证的 Gronwall 估计，得到
\[
 \sup_t\frac{|\Phi_t(a+v)-\Phi_t(a)-A(t,a)v|}{|v|}\longrightarrow0.
\]
这证明的是完整 Fréchet 微分，而非只沿坐标求导。若 \(a\) 改为邻近点，矩阵方程的系数一致收敛；作差和 Gronwall 估计又给出 \(A\) 对 \((t,a)\) 连续。因此流为 \(C^1\)。

对一段已存在的轨道，其图像在任意紧时间区间内为紧集。用上述局部构造覆盖它，并取有限子覆盖，连续依赖使邻近初始点的解也能逐段通过这些邻域。故流的定义域为开集，局部可微结论可以沿轨道继续。最后，\cref{b2:prop:07-flow-composition}给出双方均 \(C^1\) 的逆映射，所以导数矩阵可逆，像集也为开集。
\end{proof}

流的空间可微性比轨道本身可微更强。前者控制相邻粒子之间的距离、角度和体积，后者只说明单个粒子的速度存在。后面 Jacobi 行列式的演化式依赖的是\eqref{b2:eq:07-variational-ode}，不能只从每条轨道 \(C^1\) 便直接写下。

\subsection{沿特征线的表示公式}
\label{b2:subsec:070203}

对自治输运问题 \(u_t+b(x)\cdot\nabla u=0\)，沿轨道有
\[
 \frac{\mathrm d}{\mathrm dt}u\bigl(\Phi_t(a),t\bigr)=0.
\]
因此只要轨道可从时刻 \(t\) 返回 \(0\)，解就由
\begin{equation}\label{b2:eq:07-autonomous-representation}
 u(x,t)=g\bigl(\Phi_{-t}(x)\bigr)
\end{equation}
给出。这里的“返回”既包括时间存在范围，也包括轨道始终位于系数定义域；若在有界区域中轨道先碰到边界，单靠内部初值不足以继续决定数值。

\ProseParagraphBreak

例如平面速度 \(b(x,y)=(-y,x)\) 给出
\[
 \Phi_t(a_1,a_2)=(a_1\cos t-a_2\sin t,\ a_1\sin t+a_2\cos t).
\]
解是初始图形的旋转，\(u(x,y,t)=g(x\cos t+y\sin t,-x\sin t+y\cos t)\)。另一例是 \(b(x)=x\)，其流为 \(\Phi_t(a)=e^ta\)，故 \(u(x,t)=g(e^{-t}x)\)。后者虽然函数值不变，体积却按 \(e^{dt}\) 扩张；把 \(u\) 的积分解释成守恒质量就会出错。

\subsection{存在与唯一性}
\label{b2:subsec:070204}

\begin{proposition}\label{b2:prop:07-global-flow}
设 \(b\in C^1(\mathbb R^d;\mathbb R^d)\)，且存在 \(C\geq0\) 使
\( |b(x)|\leq C(1+|x|)\)。则所有轨道在整个实轴上存在。对 \(g\in C^1(\mathbb R^d)\)，式\eqref{b2:eq:07-autonomous-representation}给出整个时空上的唯一 \(C^1\) 输运解。
\end{proposition}
\begin{proof}
对正时间，积分方程与 Gronwall 估计给出
\(1+|\Phi_t(a)|\leq(1+|a|)e^{Ct}\)。负时间对反向速度 \(-b\) 使用同一估计。若最大解区间有有限右端点 \(T\)，此估计使轨道在 \([0,T)\) 落于某个紧球。该球上速度有界，所以 \(|\Phi_t(a)-\Phi_s(a)|\leq M|t-s|\)，从而 \(t\uparrow T\) 时有极限 \(a_T\)。由 \(a_T\) 出发的局部解可与原解拼接；积分方程保证拼接点导数也满足方程，与最大性矛盾。左端点相同。

全局流的 \(C^1\) 性保证表示公式给出 \(C^1\) 函数。沿特征线该函数等于 \(g(a)\)，故满足方程。任意另一经典解沿同一轨道也恒等于 \(g(a)\)，因此唯一。
\end{proof}

线性增长条件可以有许多替代，但不能不加条件地省去。对于 \(b(x)=x^2\)，从 \(a>0\) 出发的轨道为 \(a/(1-at)\)，在 \(t=1/a\) 逃向无穷。局部唯一性始终成立，障碍在于不存在覆盖所有初始点的共同正时间流。另一方面，若只关心一条有限轨道及其附近，紧集上的局部条件已经足够，不必强求整个空间上的有界导数。

% !TeX root = ../../Book2.tex
\section{变系数输运方程}
\label{b2:sec:0703}

本节把时间依赖、源项和体积变化同时纳入输运。这里 \(d\) 为正整数，\(T>0\)。为避免轨道在有限时间逃离空间，先假定 \(b\in C^1(\mathbb R^d\times[0,T];\mathbb R^d)\)，并有 \(|b(x,t)|\leq C(1+|x|)\)、\(C\geq0\)。局部版本只需把讨论限制在能够返回初始面的轨道邻域中。

\subsection{时空特征线}
\label{b2:subsec:070301}

把从时刻 \(s\) 的位置 \(a\) 出发的轨道记为 \(X(t;s,a)\)，即
\[
 \partial_tX(t;s,a)=b\bigl(X(t;s,a),t\bigr),\quad X(s;s,a)=a.
\]
扩充空间中的向量场 \(B(x,t)=(b(x,t),1)\) 是自治的。对 \(B\) 应用\cref{b2:thm:07-local-ode}与\cref{b2:thm:07-flow-derivative}，再以空间线性增长估计延拓，便得到整个 \([0,T]\) 上的轨道及其 \(C^1\) 依赖。特别地，时间坐标严格按单位速度增加，故轨道与每个等时面恰相交一次。
\symidx{\(X(t;s,a)\)：从时刻 \(s\)、位置 \(a\) 出发的轨道}

\ProseParagraphBreak

唯一性给出的复合关系现在是
\[
 X(t;r,X(r;s,a))=X(t;s,a),\quad
 X(s;t,X(t;s,a))=a.
\]
中间时刻不能省略：一般没有仅由 \(t-s\) 决定的单参数流。固定终点 \((x,t)\) 后，曲线 \(r\mapsto X(r;t,x)\) 描述该粒子的过去，初始位置为 \(X(0;t,x)\)。例如 \(b(x,t)=t x\) 时，
\[
 X(t;s,a)=\exp\left(\frac{t^2-s^2}{2}\right)a;
\]
从 \(s=0\) 走一个单位时间与从 \(s=1\) 走一个单位时间，伸缩比例不同。

\begin{proposition}[速度扰动的轨道估计]\label{b2:prop:07-speed-stability}
设 \(d\) 为正整数，\(T>0\)，\(L,\delta\geq0\)，\(b,\widetilde b\in C^1(\mathbb R^d\times[0,T];\mathbb R^d)\) 满足 \(|b(x,t)|+|\widetilde b(x,t)|\leq C(1+|x|)\)（某个 \(C\geq0\)），从 \((a,s)\)、\(a\in\mathbb R^d\) 出发的轨道分别记为 \(X,\widetilde X\)，\(s,t\in[0,T]\)，并且
\[
 \sup_{x,t}\|D_xb(x,t)\|\leq L,
 \quad\sup_{x,t}|b(x,t)-\widetilde b(x,t)|\leq\delta.
\]
则从同一时刻、同一点出发的两条轨道满足
\[
 |X(t;s,a)-\widetilde X(t;s,a)|
 \leq\delta|t-s|e^{L|t-s|}.
\]
\end{proposition}
\begin{proof}
将两条轨道方程相减，在速度差中加减 \(b(\widetilde X,r)\)。中值估计和速度差界给出，若 \(t\geq s\)，
\[
 |X(t)-\widetilde X(t)|
 \leq L\int_s^t|X(r)-\widetilde X(r)|\,\mathrm{d}r+\delta(t-s).
\]
对一个固定终点 \(t\)，把右侧最后一项在整个 \([s,t]\) 上放大为常数 \(\delta(t-s)\)，Gronwall 给出所述界。反向时间对反向速度应用相同论证，所以 \(t<s\) 也成立。
\end{proof}

若同一初值 \(g\in C^1(\mathbb R^d;\mathbb C)\) 有全空间 Lipschitz 常数 \(M\)，分别以这两个速度场求解 \(u_t+b\cdot\nabla u=0\) 与 \(\widetilde u_t+\widetilde b\cdot\nabla\widetilde u=0\)，便有 \(\|u(\cdot,t)-\widetilde u(\cdot,t)\|_\infty\leq M\delta t e^{Lt}\)、\(0\leq t\leq T\)。这里没有源项或零阶项；对终点固定的两条回溯轨道使用命题，再将它们的初始位置代入 \(g\)。例如常速度增加 \(\varepsilon\) 时，相同图形的位置误差为 \(\varepsilon t\)。若初始图形有跳跃，平移再小也可能在上确界范数中产生大小为一的误差；速度稳定性和选择何种数据范数必须同时考虑。

\subsection{源项}
\label{b2:subsec:070302}

\begin{theorem}\label{b2:thm:07-variable-transport}
设 \(d\) 为正整数，\(T>0\)，\(b\in C^1(\mathbb R^d\times[0,T];\mathbb R^d)\)，且 \(|b(x,t)|\leq C(1+|x|)\) 对某个 \(C\geq0\) 成立。记 \(X(t;s,a)\) 为 \(\partial_tX=b(X,t)\)、\(X(s;s,a)=a\) 的轨道。设 \(c,f\in C^1(\mathbb R^d\times[0,T];\mathbb C)\)、\(g\in C^1(\mathbb R^d;\mathbb C)\)。初值问题
\[
 u_t+b\cdot\nabla u+c\cdot u=f,\quad u(x,0)=g(x)
\]
存在唯一的经典解 \(u\in C^1(\mathbb R^d\times[0,T];\mathbb C)\)。记
\[
 E(t,r;x)=\exp\left(-\int_r^t c(X(q;t,x),q)\,\mathrm dq\right),
\]
则解为
\begin{equation}\label{b2:eq:07-variable-transport}
 u(x,t)=E(t,0;x)\cdot g(X(0;t,x))
 +\int_0^tE(t,r;x)\cdot f(X(r;t,x),r)\,\mathrm dr.
\end{equation}
\end{theorem}
\begin{proof}
固定 \(a\)，令 \(v(t)=u(X(t;0,a),t)\)，则链式法则把原方程变成
\[
 v'(t)+c(X(t;0,a),t)\cdot v(t)=f(X(t;0,a),t),\quad v(0)=g(a).
\]
用积分因子求解，再令 \(a=X(0;t,x)\)，即得\eqref{b2:eq:07-variable-transport}。因每条轨道的标量方程唯一，原方程也至多有一个经典解。

反过来，所给公式在紧时空集合上由 \(C^1\) 函数的复合及有限区间积分组成。轨道的连续性使相应轨道集合紧致，因此被积函数及一阶导数均有局部共同界，积分求导合法。沿 \(X(t;0,a)\) 拉回该公式时，复合关系把每条历史轨道化为 \(X(r;0,a)\)，于是恢复刚才的积分因子公式，得到方程和初值。
\end{proof}

若 \(c\) 为实值且 \(c\geq0\)，因子 \(E\) 介于 \(0\) 与 \(1\) 之间。对于非负实值初值和源项，解始终非负；对于有界复值数据，
\[
 \|u(\cdot,t)\|_\infty\leq\|g\|_\infty
 +\int_0^t\|f(\cdot,r)\|_\infty\,\mathrm dr.
\]
对一般复值 \(c\)，若存在 \(C_0\geq0\) 使 \(\operatorname{Re}c\geq-C_0\)，则 \(|E(t,r;x)|\leq e^{C_0(t-r)}\)，只需在上述上界右端补上 \(e^{C_0t}\)。这类估计还可用于两组数据之差，给出固定速度和零阶系数下的稳定性；两种上界均要求右侧的数据范数积分有限。源项按轨道积分，而零阶项以指数积累，二者对解的作用不同。

\subsection{支集传播}
\label{b2:subsec:070303}

在齐次情形 \(f=0\) 中，指数因子永不为零，因而
\begin{equation}\label{b2:eq:07-support-flow}
 \operatorname{supp}u(\cdot,t)
 =X(t;0,\operatorname{supp}g).
\end{equation}
这是闭支集的等式：流是全空间的同胚，它把非零点集合映为非零点集合，也把闭包映为闭包。

\ProseParagraphBreak

有源项时，各时刻产生的数值继续沿轨道传播，但不同来源之间可能抵消。由表示公式只能断言
\[
 \operatorname{supp}u(\cdot,t)\subseteq
 \overline{X(t;0,\operatorname{supp}g)
 \cup\bigcup_{0\leq r\leq t}X(t;r,\operatorname{supp}f(\cdot,r))}.
\]
为核对这一包含，取右侧闭集之外一点及一个与它不交的小邻域。每个邻近终点的回溯轨道都避开初始支集与全部源项支集，表示公式的每一项因此都为零。闭包不能随意去掉，因为不同时间的支集之并一般不闭。

\begin{proposition}\label{b2:prop:07-finite-speed}
设 \(d\) 为正整数，\(T>0\)，固定速度场 \(b\in C^1(\mathbb R^d\times[0,T];\mathbb R^d)\) 和零阶系数 \(c\in C^1(\mathbb R^d\times[0,T];\mathbb C)\)。设 \(M\geq0\) 且 \(|b(x,t)|\leq M\)，记 \(X(t;s,a)\) 为 \(\partial_tX=b(X,t)\)、\(X(s;s,a)=a\) 的轨道。则对 \(s,t\in[0,T]\)、\(a\in\mathbb R^d\)，有 \(|X(t;s,a)-a|\leq M|t-s|\)。

对 \(j=1,2\)，取 \(g_j\in C^1(\mathbb R^d;\mathbb C)\)、\(f_j\in C^1(\mathbb R^d\times[0,T];\mathbb C)\)，令 \(u_j\in C^1(\mathbb R^d\times[0,T];\mathbb C)\) 为同一组系数下 \((u_j)_t+b\cdot\nabla u_j+cu_j=f_j\)、\(u_j(x,0)=g_j(x)\) 的唯一经典解。对固定的 \((x,t)\in\mathbb R^d\times[0,T]\)，若 \(g_1=g_2\) 在 \(\overline B_{Mt}(x)\) 上成立，且 \(f_1=f_2\) 在回溯圆锥
\[
 \{\,(y,r)\mid0\leq r\leq t,\ |y-x|\leq M(t-r)\,\}
\]
上成立，则 \(u_1(x,t)=u_2(x,t)\)。因而在固定 \(b,c\) 后，该点的解值仅依赖于初值球内的 \(g\) 和回溯圆锥内的 \(f\)。
\end{proposition}
\begin{proof}
积分轨道方程并取绝对值得到速度界。对每个回溯时刻 \(r\)，该界给出 \(|X(r;t,x)-x|\leq M(t-r)\)，特别是初始位置属于 \(\overline B_{Mt}(x)\)。两解使用同一个流，也使用由同一 \(c\) 决定的积分因子 \(E\)。在\eqref{b2:eq:07-variable-transport}中相减，得到
\[
 \begin{aligned}
 u_1(x,t)-u_2(x,t)
 &=E(t,0;x)(g_1-g_2)(X(0;t,x))\\
 &\quad+\int_0^t E(t,r;x)(f_1-f_2)(X(r;t,x),r)\,\mathrm dr=0.
 \end{aligned}
\]
初值项和源项分别因球内和圆锥内的数据相同而为零。
\end{proof}

实际依赖集通常比这个圆锥小得多：标量输运只读取一条回溯曲线。圆锥的作用是给出不必显式求流的统一上界。线性增长的速度仍保证有限时间的轨道存在，但没有一个对所有起点通用的有限传播速度；例如 \(b(x)=x\) 在远离原点处可以运动得任意快。

\subsection{Jacobi 行列式与质量输运初步}
\label{b2:subsec:070304}

记
\[
 J(t;s,a)\coloneqq\det D_aX(t;s,a).
\]
初始时刻 \(J(s;s,a)=1\)。正、逆流的可微性保证导数矩阵一直可逆，所以行列式不能变号。下面的公式进一步说明其大小。
\symidx{\(J(t;s,a)\)：流对初始位置的 Jacobi 行列式}

\begin{proposition}[Liouville 公式]\label{b2:prop:07-jacobian}
设 \(d\) 为正整数，\(T>0\)，\(b\in C^1(\mathbb R^d\times[0,T];\mathbb R^d)\)，且 \(|b(x,t)|\leq C(1+|x|)\) 对某个 \(C\geq0\) 成立。记 \(X(t;s,a)\) 为 \(\partial_tX=b(X,t)\)、\(X(s;s,a)=a\) 的轨道。对 \(s,t\in[0,T]\)、\(a\in\mathbb R^d\)，令 \(J(t;s,a)=\det D_aX(t;s,a)\)。
\begin{equation}\label{b2:eq:07-jacobian}
 J(t;s,a)=\exp\left(\int_s^t\operatorname{div}_xb(X(r;s,a),r)\,\mathrm dr\right)>0.
\end{equation}
\end{proposition}
\begin{proof}
令 \(A=D_aX\)、\(B=D_xb(X,t)\)，则可微依赖给出 \(A'=BA\)。对可逆矩阵，行列式的多重线性展开给出
\(\det(A+H)=\det A\cdot(1+\operatorname{tr}(A^{-1}H)+o(\|H\|))\)。事实上先提出 \(A\)，再展开 \(\det(I+K)\)，所有一次项之和为 \(\operatorname{tr}K\)，其余项至少二次。因此
\[
 J'=J\cdot\operatorname{tr}(A^{-1}A')
 =J\cdot\operatorname{tr}(A^{-1}BA)
 =J\cdot\operatorname{tr}B.
\]
最后一步可由 \(\operatorname{tr}(PQ)=\sum_{i,j}p_{ij}q_{ji}=\operatorname{tr}(QP)\) 验证。解此标量方程，并使用初始值 \(J=1\)，得到结论。
\end{proof}

这就是流的\term[liouville-gongshi]{Liouville 公式}[Liouville formula]。散度描述局部体积变化率；散度为零的流保持 Lebesgue 体积。密度 \(\rho\) 要抵消体积变化，因此守恒质量的方程应写成\term[lianxuxing-fangcheng]{连续性方程}[continuity equation]
\begin{equation}\label{b2:eq:07-continuity}
 \rho_t+\operatorname{div}_x(b\rho)=0,
 \quad\text{即}\quad
 \rho_t+b\cdot\nabla\rho+(\operatorname{div}b)\rho=0.
\end{equation}

\begin{theorem}\label{b2:thm:07-mass-transport}
设 \(d\) 为正整数，\(T>0\)，\(b\in C^2(\mathbb R^d\times[0,T];\mathbb R^d)\)，\(|b(x,t)|\leq C(1+|x|)\) 对某个 \(C\geq0\) 成立；\(X(t;0,a)\) 为此速度场的轨道，\(J(t;0,a)=\det D_aX(t;0,a)\)。设 \(\rho_0\in C^1(\mathbb R^d;\mathbb C)\)。则初值问题 \(\rho_t+\operatorname{div}_x(b\rho)=0\)、\(\rho(\cdot,0)=\rho_0\) 有唯一的经典解 \(\rho\in C^1(\mathbb R^d\times[0,T];\mathbb C)\)，并满足
\begin{equation}\label{b2:eq:07-density-jacobian}
 \rho(X(t;0,a),t)\cdot J(t;0,a)=\rho_0(a).
\end{equation}
若 \(\rho_0\) 为非负实值函数，或 \(\rho_0\in L^1(\mathbb R^d;\mathbb C)\)，则对每个 Borel 集 \(E\subseteq\mathbb R^d\) 和 \(t\in[0,T]\)，
\[
 \int_{X(t;0,E)}\rho(x,t)\,\mathrm dx=\int_E\rho_0(a)\,\mathrm da.
\]
非负情形允许两边为 \(+\infty\)。可积情形还有 \(\|\rho(\cdot,t)\|_1=\|\rho_0\|_1\)。
\end{theorem}
\begin{proof}
因为 \(\operatorname{div}b\in C^1\)，可对\eqref{b2:eq:07-continuity}使用\cref{b2:thm:07-variable-transport}，得到经典解的存在唯一性。沿轨道，\(\rho\) 满足的标量方程与 \(J\) 的方程恰有相反的零阶系数，乘积求导为零，得到\eqref{b2:eq:07-density-jacobian}。对微分同胚 \(a\mapsto X(t;0,a)\) 使用\secref{b2:sec:0102}的换元公式，因为 \(J>0\)，得到积分等式。若初始密度非负，可直接用非负换元；若绝对可积，先对绝对值换元，得到可积性及范数恒等式，再对正、负部或实、虚部分别积分。
\end{proof}

这里加 \(C^2\) 条件是为了由显式公式得到 \(C^1\) 密度；\cref{b2:prop:07-jacobian}本身只需要 \(b\in C^1\)。对于速度 \(b(x)=x\)，输运标量为 \(u(x,t)=g(e^{-t}x)\)，而守恒密度为 \(\rho(x,t)=e^{-dt}\cdot\rho_0(e^{-t}x)\)。两者的支集作同样的扩张，后者的振幅同时降低。运动几何、所携带的值与密度三者至此有了各自准确的表达式。

\ProseParagraphBreak

\Cref{b2:tab:07-scalar-density}把没有源项时的两种未知量并列比较。表中 \(d\) 为正整数，\(T>0\)，同一实速度场 \(b\in C^2(\mathbb R^d\times[0,T];\mathbb R^d)\) 满足 \(|b(x,t)|\leq C(1+|x|)\)、\(C\geq0\)；记 \(X(t,a)=X(t;0,a)\)、\(J(t,a)=\det D_aX(t;0,a)>0\)。标量初值取 \(g\in C^1(\mathbb R^d;\mathbb C)\)，密度初值取非负的 \(\rho_0\in C^1(\mathbb R^d;\mathbb R)\cap L^1(\mathbb R^d;\mathbb R)\)。表中的“移动区域”是同一群初始粒子 \(E\) 的像 \(X(t,E)\)，其中 \(E\subseteq\mathbb R^d\) 为 Borel 集。

\begin{table}[htbp]
\centering
\small
\renewcommand{\arraystretch}{1.3}
\caption[输运标量与守恒密度]{输运标量与守恒密度。两列使用同一个流；密度沿轨道的变化抵消了体积因子，标量沿轨道的数值不变则没有这种补偿。}
\label{b2:tab:07-scalar-density}
\begin{tabularx}{\linewidth}{@{}>{\raggedright\arraybackslash}p{1.5cm}>{\raggedright\arraybackslash}X>{\raggedright\arraybackslash}X@{}}
\toprule
 & 输运标量 \(u\) & 守恒密度 \(\rho\) \\
\midrule
方程 & \(u_t+b\cdot\nabla u=0\) & \(\rho_t+\operatorname{div}_x(b\rho)=0\) \\
沿轨道 & \(u(X(t,a),t)=g(a)\) & \(\rho(X(t,a),t)J(t,a)=\rho_0(a)\) \\
\(b(x)=x\) & \(u(x,t)=g(e^{-t}x)\) & \(\rho(x,t)=e^{-dt}\rho_0(e^{-t}x)\) \\
保持的量 & 每个粒子携带的函数值；总积分或 \(L^p\) 范数即使有限，也不必保持。 & 随流移动区域内的质量：\(\int_{X(t,E)}\rho=\int_E\rho_0\)。总质量及 \(L^1\) 范数不变。 \\
\bottomrule
\end{tabularx}
\end{table}

移动区域的质量不变，并不表示每个固定观察区域内的质量也不变；粒子仍可穿过固定区域的边界。这一区别不要求移动区域具有光滑边界，质量恒等式已经由换元对所有 Borel 集给出。

\ProseParagraphBreak

Jacobi 公式同时决定各个积分范数的演化。以下取正整数 \(d\)、\(T>0\)，实速度场 \(b\in C^1(\mathbb R^d\times[0,T];\mathbb R^d)\) 满足 \(|b(x,t)|\leq C(1+|x|)\)、\(C\geq0\)，且 \(c\in C^1(\mathbb R^d\times[0,T];\mathbb R)\)、\(g\in C^1(\mathbb R^d;\mathbb C)\)。对无源问题 \(u_t+b\cdot\nabla u+cu=0\)、\(u(x,0)=g(x)\) 的经典解，沿轨道有
\[
 u(X(t;0,a),t)=g(a)\exp\left(-\int_0^tc(X(r;0,a),r)\,\mathrm{d}r\right).
\]
对 \(1\leq p<\infty\) 作换元，得到精确式
\[
 \|u(\cdot,t)\|_p^p
 =\int |g(a)|^p\exp\left(\int_0^t
   (\operatorname{div}b-pc)(X(r;0,a),r)\,\mathrm{d}r\right)\,\mathrm{d}a.
\]
若存在实常数 \(D,c_0\)，使整个时空条带上 \(\operatorname{div}b\leq D\)、\(c\geq c_0\)，且 \(g\in L^p(\mathbb R^d;\mathbb C)\)，便有 \(\|u(\cdot,t)\|_p\leq e^{(D/p-c_0)t}\|g\|_p\)。无穷范数不读取体积；若 \(g\in L^\infty\)，直接由表示公式得 \(\|u(\cdot,t)\|_\infty\leq e^{-c_0t}\|g\|_\infty\)。保持同一组 \(b,c\)，在有源情形取 \(f\in C^1(\mathbb R^d\times[0,T];\mathbb C)\)，对每个源时刻用同一换元，再用 Minkowski 不等式，给出
\[
 \|u(\cdot,t)\|_p
 \leq e^{(D/p-c_0)t}\|g\|_p
 +\int_0^t e^{(D/p-c_0)(t-r)}\|f(\cdot,r)\|_p\,\mathrm{d}r,
\]
其中 \(p=\infty\) 按 \(D/p=0\) 解释，并要求右侧有限。

\ProseParagraphBreak

对\cref{b2:thm:07-mass-transport}中 \(b\in C^2\) 所给出的守恒密度，零阶系数是 \(c=\operatorname{div}b\)。它的 \(p\) 次质量因子变为 \((1-p)\operatorname{div}b\)。所以只有 \(p=1\) 对任意体积变化恰好抵消；若 \(1\leq p\leq\infty\)、\(\rho_0\in L^p(\mathbb R^d;\mathbb C)\)，且整个时空条带上 \(|\operatorname{div}b|\leq D\)、\(D\geq0\)，则
\[
 \|\rho(\cdot,t)\|_p\leq e^{(1-1/p)Dt}\|\rho_0\|_p.
\]
对于 \(b(x)=x\)，还可直接算出 \(\|\rho(\cdot,t)\|_p=e^{-d(1-1/p)t}\|\rho_0\|_p\)，而输运标量的范数为 \(e^{dt/p}\|g\|_p\)。同一流对不同物理量产生不同范数变化，不能把所有输运都当成保持每个 \(L^p\) 范数的平移。

% !TeX root = ../../Book2.tex
\OptionalSection{一阶非线性方程}{b2:sec:0704}

在线性输运中，先由已知速度场确定轨道，再沿轨道计算数值。若速度依赖未知函数，这两步便相互耦合。本节只构造局部光滑解，并用一个显式例子说明这种解可能在有限时间失效。后面的波动、热与 Laplace 方程不依赖本节。

\subsection{拟线性方程}
\label{b2:subsec:070401}

考虑已在\secref{b2:sec:0104}定义的拟线性类型
\[
 u_t+b(x,t,u)\cdot\nabla u=f(x,t,u),\quad u(x,0)=g(x).
\]
设 \(b,f\) 在相关变量的开集上为 \(C^1\)，\(g\in C^1\)。把未知的函数值也看成一个坐标，解联立系统
\begin{equation}\label{b2:eq:07-quasilinear-system}
 \begin{aligned}
 X'&=b(X,t,Z),&X(0,a)&=a,\\
 Z'&=f(X,t,Z),&Z(0,a)&=g(a).
 \end{aligned}
\end{equation}
对向量 \((X,t,Z)\) 应用局部流定理，得到其对 \((t,a)\) 的 \(C^1\) 依赖。时刻零有 \(D_aX=I\)，故映射 \((a,t)\mapsto(X(t,a),t)\) 在每个 \((a_0,0)\) 附近的导数可逆。

\begin{proposition}\label{b2:prop:07-quasilinear-local}
在上述假设下，对每个初始点 \(a_0\)，存在它在时空中的邻域，在其中拟线性初值问题有唯一的局部 \(C^1\) 解。
\end{proposition}
\begin{proof}
由\cref{b2:thm:01-inverse}，缩小邻域后可反解 \(a=A(x,t)\)。令 \(u(x,t)=Z(t,A(x,t))\)，则
\(u(X(t,a),t)=Z(t,a)\)。对此式求时间导数并代入\eqref{b2:eq:07-quasilinear-system}，得到偏微分方程。初始时 \(A(x,0)=x\)，所以 \(u(x,0)=g(x)\)。若另有经典解 \(v\)，沿满足 \(X'=b(X,t,v(X,t))\) 的曲线令 \(Z=v(X,t)\)，链式法则表明 \((X,Z)\) 满足同一联立系统和同一初值。常微分方程唯一性使二者与构造的 \((X,Z)\) 相同，再由局部反演得到 \(v=u\)。
\end{proof}

这一结论的邻域可以依赖 \(a_0\)。即使联立系统一直有解，其空间投影 \(a\mapsto X(t,a)\) 也可能先失去可逆性。下面的例子把这种障碍显示为一个分母趋零的计算。

\subsection{特征线交叉}
\label{b2:subsec:070402}

\term[burgers-fangcheng]{无黏 Burgers 方程}[inviscid Burgers equation]为
\[
 u_t+u u_x=0.
\]
沿特征线 \(Z'=0\)，所以 \(Z(t,a)=g(a)\)、\(X(t,a)=a+tg(a)\)。若 \(g\in C^1(\mathbb R)\) 有界，且 \(1+tg'(a)>0\) 对所有 \(a\) 成立，则 \(X(t,\cdot)\) 严格递增，又因 \(|X(t,a)-a|\leq t\|g\|_\infty\) 而满射。它是全实轴的微分同胚，解由 \(u(X(t,a),t)=g(a)\) 给出，并满足
\begin{equation}\label{b2:eq:07-burgers-gradient}
 u_x(X(t,a),t)=\frac{g'(a)}{1+tg'(a)}.
\end{equation}
若 \(g'\geq0\)，上述构造对所有正时间有效；若 \(g'\) 的负最小值 \(m\) 在某点取到，则在 \(T_*=-1/m\) 之前有效，而该点轨道上的梯度在 \(t\uparrow T_*\) 时趋于负无穷。

\ProseParagraphBreak

具体取 \(g(a)=-\tanh a\)，有 \(m=-1\)、\(T_*=1\)。函数值始终处于 \((-1,1)\)，但原点处梯度为 \(-1/(1-t)\)。失效不是振幅无限增长，而是相邻初始点的空间间隔被压缩为零。当 \(t>1\) 时，映射 \(a-t\tanh a\) 在原点附近下降，在两端上升，因而不单射。两条不同初始标签的轨道可以投影到同一个时空点，而所携带的值不同，不可能成为单值经典解。\Cref{b2:fig:07-characteristic-compression}区分了光滑阶段与失效后的轨道投影。
% 自绘：g(a)=-tanh(a) 的特征投影；虚线仅表示光滑失效后的轨道投影。
\begin{figure}[htbp]
\centering
\begin{tikzpicture}[x=2.3cm,y=2.2cm,>=stealth]
 \draw[->] (-1.35,0)--(1.35,0) node[right]{\(x\)};
 \draw[->] (-1.3,-0.05)--(-1.3,1.6) node[above]{\(t\)};
 \draw[gray,densely dotted] (-1.25,1)--(1.25,1);
 \node[left] at (-1.3,1){\(1\)};
 \foreach \a/\v in {-1/-0.761594,-0.75/-0.635149,-0.5/-0.462117,-0.25/-0.244919,0/0,0.25/0.244919,0.5/0.462117,0.75/0.635149,1/0.761594}{
   \draw[AnalysisBlue,thick] (\a,0)--({\a-\v},1);
   \draw[AnalysisBlue,dashed] ({\a-\v},1)--({\a-1.45*\v},1.45);
 }
 \foreach \x in {-1,0,1}{\draw (\x,0.025)--(\x,-0.025) node[below]{\(\x\)};}
 \node[anchor=west,align=left] at (0.68,1.28){\small 轨道投影};
 \node[anchor=west] at (0.68,0.72){\small 光滑阶段};
\end{tikzpicture}
\caption[特征线的压缩与投影交叉]{特征线的压缩与投影交叉。初值为 \(g(a)=-\tanh a\)，轨道投影为 \(x=a-t\tanh a\)。在 \(t=1\) 首次发生梯度爆破；虚线延长只展示投影的交叉，不表示单值经典解的延续。}
\label{b2:fig:07-characteristic-compression}
\end{figure}


\subsection{激波形成的预告}
\label{b2:subsec:070403}

在经典解存在时，Burgers 方程等价于
\(u_t+\partial_x(u^2/2)=0\)。对于有跳跃的函数，后一种写法仍可通过分布导数解释。若一般通量 \(F\in C^1(\mathbb R)\)，一条 \(C^1\) 曲线 \(x=s(t)\) 两侧的函数都是常数 \(u_-,u_+\)，把测试积分在曲线两侧分别分部积分，界面项为
\[
 \int \bigl(F(u_+)-F(u_-)-s'(t)(u_+-u_-)\bigr)
 \varphi(s(t),t)\,\mathrm dt.
\]
它对所有紧支集光滑测试函数为零，当且仅当
\[
 s'(t)(u_+-u_-)=F(u_+)-F(u_-).
\]
这称为\term[rankine-hugoniot]{Rankine--Hugoniot 条件}[Rankine--Hugoniot condition]。对不同的 Burgers 左右状态，速度为 \((u_-+u_+)/2\)；这样的移动跳跃称为\term[jibo]{激波}[shock wave]。

\ProseParagraphBreak

界面恒等式只给出弱方程的相容条件，没有证明某个跳跃是光滑解失效后的唯一延续。物理上允许的压缩跳跃与不允许的展开跳跃，都可能满足同一恒等式。要挑选延续还需附加条件，并建立存在、稳定性与唯一性理论；本节的梯度爆破计算本身不承担这些结论。

\subsection{Hamilton--Jacobi 方程}
\label{b2:subsec:070404}

另一类重要的一阶非线性方程是\term[hamilton-jacobi-fangcheng]{Hamilton--Jacobi 方程}[Hamilton--Jacobi equation]
\[
 u_t+H(x,\nabla u)=0.
\]
这里 \(H\) 为给定函数。先假定 \(H\in C^3\)、\(g\in C^3\)，只在初始点附近构造光滑解。取
\begin{equation}\label{b2:eq:07-hamilton-system}
 X'=H_p(X,P),\quad P'=-H_x(X,P),\quad
 Z'=P\cdot H_p(X,P)-H(X,P),
\end{equation}
初始值为 \(X(0,a)=a\)、\(P(0,a)=\nabla g(a)\)、\(Z(0,a)=g(a)\)。前两式描述位置和候选梯度的演化，第三式计算函数值。

\ProseParagraphBreak

这里需要的二阶参数可微性仍来自流构造：对 \(C^2\) 向量场 \(B\)，在\eqref{b2:eq:07-variational-ode}上再取差商，二阶候选导数满足
\[
 V'=DB(Y)V+D^2B(Y)[A v,A w],\quad V(0)=0.
\]
右侧系数连续，矩阵积分方程给出唯一 \(V\)。以 \(D^2B\) 在紧集上的一致连续性控制差商余项，再用 Gronwall 估计，得到余项一致趋零。这与一阶证明的线性余项论证相同，且明确给出二阶导数的方程；所以此处构造的 \(X,P,Z\) 对 \((t,a)\) 为 \(C^2\)。

\begin{proposition}\label{b2:prop:07-hamilton-local}
在上述假设下，系统\eqref{b2:eq:07-hamilton-system}在充分小的邻域中确定一个 \(C^2\) 解 \(u\)，满足 Hamilton--Jacobi 方程及初值 \(g\)。
\end{proposition}
\begin{proof}
从 \(D_aX(0,a)=I\) 出发，局部反演得到 \(a=A(x,t)\)，令 \(u(x,t)=Z(t,A(x,t))\)。关键是验证 \(P=\nabla u\)，不能把候选梯度直接当成真实梯度。对每个初始坐标 \(a_j\)，记
\(W_j=\partial_{a_j}Z-P\cdot\partial_{a_j}X\)。对系统求导，利用 \(X'=H_p\)、\(P'=-H_x\)，得到
\[
 \begin{aligned}
 \partial_{a_j}Z'
 &=P\cdot\partial_{a_j}(H_p)-H_x\cdot\partial_{a_j}X,\\
 (P\cdot\partial_{a_j}X)'
 &=-H_x\cdot\partial_{a_j}X+P\cdot\partial_{a_j}(H_p).
 \end{aligned}
\]
故 \(W_j'=0\)，初值又为 \(\partial_jg-\partial_jg=0\)。因此 \(D_aZ=P^TD_aX\)，由 \(D_aX\) 可逆得到 \(\nabla_xu(X,t)=P\)。最后沿轨道有
\(u_t+\nabla u\cdot X'=Z'=P\cdot H_p-H\)，约去相同项得 \(u_t+H(x,\nabla u)=0\)。
\end{proof}

取 \(H(p)=|p|^2/2\) 时，\(P=\nabla g(a)\)、\(X=a+t\nabla g(a)\)、\(Z=g(a)+t|\nabla g(a)|^2/2\)。可逆性要求 \(I+tD^2g(a)\) 不退化。这与 Burgers 方程的分母具有同一几何含义，但高维中须检查整个矩阵。局部光滑构造失效之后的解概念，本章不再展开。

\subsection{守恒律的位置}
\label{b2:subsec:070405}

形如 \(u_t+\operatorname{div}_xF(u)=0\) 的方程称为\term[shouhenglv]{守恒律}[conservation law]，其中 \(F:\mathbb R\rightarrow\mathbb R^d\)。光滑时，它等价于拟线性方程 \(u_t+F'(u)\cdot\nabla u=0\)；不光滑时，两种写法的意义不同，后一项中的乘积未必已有定义。守恒形式可直接作用于测试函数，通过积分表达局部通量平衡。

\ProseParagraphBreak

线性连续性方程的速度由外部给定，密度与体积的补偿由 Jacobi 公式控制；非线性守恒律的速度随状态变化，投影交叉会使经典表示失效。两者都使用特征线，却不能把线性流的全局微分同胚结论原样移到非线性问题。本节建立了光滑阶段的局部构造和一个明确的失效机制，为以后阅读守恒律与 Hamilton--Jacobi 理论划出入口。


\begin{chaptersummary}
Schwartz 空间用所有阶导数和所有多项式权重控制函数。Fourier 变换在其中交换微分与乘法，反演使它成为连续自同构；Hermite 选读专题则以正交基把这一变换的四种特征值具体呈现出来。温和分布通过作用于测试函数定义，双线性配对把已有的变换公式连续延伸到 Dirac 测度、多项式以及全部 \(L^p\) 函数。

\ProseParagraphBreak

常系数方程的频率表达是多项式乘法，但零点处不能随意相除。基本解的构造必须兼顾局部可积性、原点质量和无穷远增长。负 Laplace 的幂核、对数核和一维线性核各反映一种维数边界；与紧支集光滑右端卷积可解全空间方程，却不自动决定边值条件。下一章将用同一测试空间把离散格点与连续频率联系起来。
\end{chaptersummary}

\BookExercises

以下题目把半范数、差商和频率符号用于具体对象。前几题强调收敛与符号，后几题处理奇性、频率特征结构和方程。\Cref{b2:ex:05-hermite-oscillator}对应选读的 Hermite 专题；\cref{b2:ex:05-polynomial-comb}比较离散测度的多项式增长与指数增长，最后的\cref{b2:ex:05-cutoff-laplace-kernel}考察截断基本解后出现的壳层余项。

% 自拟：所有导数一致趋零仍不控制多项式权重。
% 预留题头及首个公式所需空间，避免题头换页后超链接锚点仍在前页。
\Needspace{6\baselineskip}
\begin{exercise}\label{b2:ex:05-schwartz-convergence}
设非零 \(\rho\in C_c^\infty(\mathbb R;\mathbb C)\)，对正整数 \(n\) 令 \(\varphi_n(x)=n^{-1}\rho(x/n)\)。证明每个固定阶导数都在 \(\mathbb R\) 上一致趋于零，但 \(\varphi_n\) 不在 \(\mathcal S(\mathbb R;\mathbb C)\) 中趋零。它作为温和分布是否逐测试函数配对趋于零？
\end{exercise}
\begin{solution}
有 \(\|\varphi_n^{(k)}\|_\infty=n^{-k-1}\|\rho^{(k)}\|_\infty\to0\)。选 \(x_0\ne0\) 使 \(\rho(x_0)\ne0\)，则 \(p_1(\varphi_n)\geq n^{-1}(1+n|x_0|)|\rho(x_0)|\)，下极限为正，故不按 Schwartz 半范数趋零。另一方面，对 \(\psi\in\mathcal S\)，\(|\int\varphi_n\psi|\leq n^{-1}\|\rho\|_\infty\cdot\|\psi\|_1\to0\)，所以在 \(\mathcal S'\) 中趋零。
\end{solution}

% 自拟：差分点质量的极限与傅里叶符号同时核对。
\Needspace{6\baselineskip}
\begin{exercise}\label{b2:ex:05-delta-differences}
对正整数 \(n\)，在 \(\mathcal S'(\mathbb R;\mathbb C)\) 中按逐测试函数配对收敛求下列分布列的极限，并计算极限的 Fourier 变换，其中 \(\delta_a\) 表示位于 \(a\in\mathbb R\) 的 Dirac 测度：
\[
 n(\delta_{1/n}-\delta_0),\quad
 n^2(\delta_{1/n}+\delta_{-1/n}-2\delta_0).
\]
\end{exercise}
\begin{solution}
对 \(\varphi\in\mathcal S\) 使用一阶与二阶 Taylor 差商，配对极限分别为 \(\varphi'(0)\) 和 \(\varphi''(0)\)，故极限为 \(-\delta_0'\)、\(\delta_0''\)。它们的变换分别为 \(-2\pi\mathrm{i}\xi\)、\(-4\pi^2\xi^2\)。也可从指数差商核对：\(n(e^{-2\pi\mathrm{i}\xi/n}-1)\to-2\pi\mathrm{i}\xi\)，其模由 \(2\pi|\xi|\) 控制；对称二阶差商的模由 \(4\pi^2\xi^2\) 控制，故两者都可对 Schwartz 测试函数用控制收敛。
\end{solution}

% 自拟：乘以x的核及原点导数项，训练分布除法。
\Needspace{6\baselineskip}
\begin{exercise}\label{b2:ex:05-x-times-distribution}
设 \(T\in\mathcal S'(\mathbb R;\mathbb C)\)。证明 \(xT=0\) 当且仅当存在 \(c\in\mathbb C\) 使 \(T=c\delta_0\)。再求方程 \(xT=\delta_0\) 在同一空间中的全部解。
\end{exercise}
\begin{solution}
取 \(\rho\in C_c^\infty\)，在原点附近为 \(1\)。任意 \(\varphi\in\mathcal S\) 可写成 \(\varphi=\varphi(0)\rho+x\psi\)，其中 \(\psi=(\varphi-\varphi(0)\rho)/x\) 在原点用光滑延拓，远处除以 \(x\) 保持快速下降，所以 \(\psi\in\mathcal S\)。若 \(xT=0\)，便有 \(\langle T,\varphi\rangle=\varphi(0)\langle T,\rho\rangle\)，即 \(T=c\delta_0\)。反向直接成立。又因 \(\langle x\delta_0',\varphi\rangle=-(x\varphi)'(0)=-\varphi(0)\)，一个特解为 \(-\delta_0'\)。全部解因而为 \(-\delta_0'+c\delta_0\)。
\end{solution}

% 自拟：主值分布、奇偶性和一阶分布方程组合得到变换常数。
\Needspace{6\baselineskip}
\begin{optionalexercise}\label{b2:ex:05-principal-value}
证明对每个 \(\varphi\in\mathcal S(\mathbb R;\mathbb C)\)，公式
\[
 \langle T,\varphi\rangle
 =\lim_{\varepsilon\downarrow0}\int_{|x|>\varepsilon}\dfrac{\varphi(x)}x\,\mathrm{d}x
\]
定义 \(T\in\mathcal S'(\mathbb R;\mathbb C)\)，称为 \(\operatorname{pv}(1/x)\)。证明 \(xT=1\)，其中右侧是常函数 \(1\) 给出的分布，并计算 \(\widehat T\)。
\end{optionalexercise}
\begin{solution}
在 \(|x|<1\) 内减去 \(\varphi(0)\)，常数项因对称性抵消，剩余绝对值由 \(\|\varphi'\|_\infty\) 控制；在 \(|x|\geq1\) 内用 \(p_2(\varphi)(1+|x|)^{-2}/|x|\) 控制。因此极限存在且满足有限半范数界。代入 \(x\varphi\) 得 \(xT=1\)，并可换元看出 \(T\) 为奇分布。

由变换公式，\(\frac{\mathrm{i}}{2\pi}\partial_\xi\widehat T=\delta_0\)，故 \(\partial_\xi\widehat T=-2\pi\mathrm{i}\delta_0\)。函数 \(-\pi\mathrm{i}\operatorname{sgn}\xi\) 的分布导数正好为右端，所以两者之差的导数为零。导数为零的温和分布是常数：对积分为零的 \(\psi\in\mathcal S\)，其原函数 \(\int_{-\infty}^x\psi\) 属于 \(\mathcal S\)，故该分布在所有零积分测试函数上为零；再减去一个固定单位积分测试函数的倍数即可。于是 \(\widehat T=-\pi\mathrm{i}\operatorname{sgn}\xi+c\)。Fourier 变换保持奇偶性，而 \(T\) 为奇，所以 \(c=0\)。
\end{solution}

% 自拟：求Hermite低阶式，并把频率特征关系与谐振子联系。
\Needspace{6\baselineskip}
\begin{optionalexercise}\label{b2:ex:05-hermite-oscillator}
在 \(\mathbb R\) 上令 \(h_0(x)=2^{1/4}e^{-\pi x^2}\)、\(a^*=(4\pi)^{-1/2}(-\partial_x+2\pi x)\)，并对 \(n\in\mathbb N_0\) 令 \(h_n=(a^*)^nh_0/\sqrt{n!}\)，即采用\subsecref{b2:subsec:050204}的归一化。求 \(h_1,h_2\) 的显式公式，证明对每个 \(n\in\mathbb N_0\) 有
\[
 \left(-\dfrac{\mathrm{d}^2}{\mathrm{d}x^2}+4\pi^2x^2\right)h_n
 =(4n+2)\pi h_n,
\]
并把 \(x^2h_0\) 展成 Hermite 函数的有限线性组合。
\end{optionalexercise}
\begin{solution}
从上升算子计算得 \(h_1=\sqrt{4\pi}\,xh_0\)、\(h_2=2^{-1/2}(4\pi x^2-1)h_0\)。记下降算子 \(a=(4\pi)^{-1/2}(\partial_x+2\pi x)\)，展开算子有 \(4\pi a^*a=-\mathrm{d}^2/\mathrm{d}x^2+4\pi^2x^2-2\pi\)。由 \(a^*ah_n=nh_n\) 得到特征方程。第二个显式公式给出 \(x^2h_0=(h_0+\sqrt2h_2)/(4\pi)\)。其变换为 \((h_0-\sqrt2h_2)/(4\pi)\)，也可与两次频率求导公式核对。
\end{solution}

% 自拟：有界光滑乘子若导数失控，可以破坏S。
\Needspace{6\baselineskip}
\begin{exercise}\label{b2:ex:05-fast-multiplier}
在 \(\mathbb R\) 上令 \(a(x)=e^{\mathrm{i}e^{x^2}}\)、\(\varphi(x)=e^{-x^2}\)。证明 \(a\colon\mathbb R\to\mathbb C\) 光滑且有界，\(\varphi\in\mathcal S(\mathbb R;\mathbb C)\)，但 \(a\varphi\notin\mathcal S(\mathbb R;\mathbb C)\)。再判断乘以 \(b(x)=e^{\mathrm{i}x^3}\) 是否给出 \(\mathcal S(\mathbb R;\mathbb C)\) 的连续自同构，并说明其逆。
\end{exercise}
\begin{solution}
直接求导得 \((a\varphi)'=a(-2xe^{-x^2}+2\mathrm{i}x)\)，其模随 \(|x|\) 无界，故不满足 Schwartz 条件。对 \(b\)，每个导数都是一个多项式乘 \(b\)，因而至多多项式增长。由正文的乘法命题，乘以 \(b\) 连续保持 \(\mathcal S\)；其逆乘子 \(\overline b\) 也有相同性质，所以还给出连续自同构。
\end{solution}

% 自拟：尺度决定屏蔽算子的基本解与L1质量。
\Needspace{6\baselineskip}
\begin{exercise}\label{b2:ex:05-resolvent-scale}
设 \(d\) 为正整数，\(E\in L^1(\mathbb R^d;\mathbb R)\) 为几乎处处由
\[
 E(x)=\int_0^\infty e^{-t}(4\pi t)^{-d/2}e^{-|x|^2/(4t)}\,\mathrm dt
\]
给出的非负等价类；允许这个积分在原点取无穷值。对 \(a>0\)，在几乎处处意义下令 \(E_a(x)=a^{d/2-1}E(\sqrt a\,x)\)。证明 \(E_a\) 是 \(a-\Delta\) 的可积基本解，求其变换和 \(L^1(\mathbb R^d)\) 范数。若 \(1\leq p\leq\infty\)、\(f\in L^p(\mathbb R^d;\mathbb C)\)，证明 \(\|E_a*f\|_p\leq a^{-1}\|f\|_p\)。
\end{exercise}
\begin{solution}
\Cref{b2:prop:05-bessel-fundamental}已证明 \(E\) 的质量为 \(1\)，且 \(\widehat E(\xi)=(1+4\pi^2|\xi|^2)^{-1}\)。尺度公式给出
\[
 \widehat E_a(\xi)=a^{-1}\widehat E(\xi/\sqrt a)
 =\dfrac1{a+4\pi^2|\xi|^2}.
\]
乘以符号 \(a+4\pi^2|\xi|^2\) 后得到 \(1\)，故为基本解。它非负，换元得 \(\|E_a\|_1=a^{-1}\|E\|_1=a^{-1}\)。最后用 Young 卷积不等式。这个估计允许任意 \(1\leq p\leq\infty\)，但不单凭此范数界宣称解具有任意阶强导数。
\end{solution}

% 自拟：离散测度的温和性及其不可微函数式解释。
\Needspace{6\baselineskip}
\begin{exercise}\label{b2:ex:05-polynomial-comb}
设 \(d\) 为正整数，\(c_k\in\mathbb C\)（\(k\in\mathbb Z^d\)），并存在与 \(k\) 无关的常数 \(C,M\geq0\)，使 \(|c_k|\leq C(1+|k|)^M\)。证明 \(T=\sum_{k\in\mathbb Z^d}c_k\delta_k\) 由绝对收敛配对定义温和分布，且截断 \(T_N=\sum_{|k|_\infty\leq N}c_k\delta_k\) 在 \(\mathcal S'(\mathbb R^d;\mathbb C)\) 中逐测试函数配对趋向 \(T\)，其中 \(N\) 遍历正整数，\(|k|_\infty=\max_j|k_j|\)。再在 \(d=1\) 时证明正测度
\[
 \mu=\sum_{k\in\mathbb Z}e^{|k|}\delta_k
\]
虽给出紧支集测试函数上的普通分布，却不能延拓为 \(\mathcal S'(\mathbb R;\mathbb C)\) 中的元素。可选取非负 \(\rho\in C_c^\infty(\mathbb R;\mathbb R)\)，使 \(\operatorname{supp}\rho\subset(-1/2,1/2)\)、\(\rho(0)=1\)，检查 \(\varphi_n(x)=e^{-n/2}\rho(x-n)\) 的全部 Schwartz 半范数与 \(\int\varphi_n\,\mathrm d\mu\)。
\end{exercise}
\begin{solution}
选整数 \(N>M+d\)。则
\[
 \sum_k|c_k\varphi(k)|
 \leq C p_N(\varphi)\sum_k(1+|k|)^{M-N}<\infty.
\]
最后的格点级数可按整数半径或二进环带比较体积：半径约 \(2^j\) 的点数至多为 \(C_d2^{jd}\)，所以所得几何级数收敛。这同时给出有限半范数界。对每个固定测试函数，绝对收敛级数的截断尾部趋零，正好是 \(\mathcal S'\) 收敛。

指数权测度在每个紧集上只有有限个格点，故局部有限；对有共同紧支集的测试函数，积分由该紧集上的有限质量乘函数的上确界控制，所以它确实定义普通分布。取题中 \(\rho\) 和正整数 \(n\)。对每个固定非负整数 \(m\)，\(\varphi_n\) 的支集在 \((n-1/2,n+1/2)\) 内，其各阶导数都是 \(e^{-n/2}\) 乘相应平移导数。因此
\[
 p_m(\varphi_n)
 \leq C_{m,\rho}\cdot e^{-n/2}(1+n)^m\longrightarrow0.
\]
故 \(\varphi_n\to0\) 于 \(\mathcal S(\mathbb R;\mathbb C)\)。另一方面，支集中只有整数 \(n\)，所以
\[
 \int\varphi_n\,\mathrm d\mu
 =e^n\cdot e^{-n/2}\rho(0)=e^{n/2}\longrightarrow\infty.
\]
若存在延拓到 \(\mathcal S\) 的连续泛函，它在这些紧支集测试函数上的值必须保持上述积分值，却又应由连续性趋于零，矛盾。这不只是说明某个增长判别不能套用，而是直接排除了连续延拓。
\end{solution}

% 自拟：移动Laplace极点并截断空间，求平滑壳层余项及其总质量；不重复正文未截断核的验证题。
\Needspace{6\baselineskip}
\begin{exercise}\label{b2:ex:05-cutoff-laplace-kernel}
设 \(d\in\{2,3\}\)、\(a\in\mathbb R^d\)、\(R>0\)。对 \(x\ne a\)，令
\[
 \Phi_{d,a}(x)=
 \begin{cases}
 -\dfrac1{2\pi}\log|x-a|,&d=2,\\[4pt]
 \dfrac1{4\pi|x-a|},&d=3.
 \end{cases}
\]
在 \(x=a\) 处任意赋一个有限值。取实值 \(\chi\in C_c^\infty(\mathbb R^d;\mathbb R)\)，使 \(\chi=1\) 于 \(B(0,1)\)、\(\chi=0\) 于 \(\{|x|\geq2\}\)，并令 \(\chi_R(x)=\chi((x-a)/R)\)、\(K(x)=\chi_R(x)\cdot\Phi_{d,a}(x)\)。
\begin{enumerate}
\item\label{b2:item:05-shifted-harmonic-flux} 证明 \(\Phi_{d,a}\) 在 \(\mathbb R^d\setminus\{a\}\) 上调和，且对每个 \(\varepsilon>0\)，
\[
 -\int_{\partial B(a,\varepsilon)}\partial_n\Phi_{d,a}\,\mathrm dS=1,
\]
其中 \(n\) 是小球的外单位法向。说明删去小球后，剩余区域在内边界处的外法向为何取相反方向。
\item\label{b2:item:05-cutoff-shell-remainder} 证明 \(K\in L^1(\mathbb R^d;\mathbb R)\) 且紧支集。对 \(x\ne a\) 令
\[
 q(x)=-2\nabla\chi_R(x)\cdot\nabla\Phi_{d,a}(x)
       -(\Delta\chi_R(x))\cdot\Phi_{d,a}(x).
\]
证明 \(q\) 在 \(a\) 附近补成零后属于 \(C_c^\infty(\mathbb R^d;\mathbb R)\)，支集包含于 \(\{\,x\mid R\leq|x-a|\leq2R\,\}\)，并通过删球积分和极限证明
\[
 -\Delta K=\delta_a+q\quad\text{于 }\mathcal S'(\mathbb R^d;\mathbb C).
\]
\item\label{b2:item:05-shell-mass} 证明 \(\int_{\mathbb R^d}q=-1\)。解释为何把基本解截成紧支集后，不能仍只留下一个点质量而没有额外余项。
\end{enumerate}
\end{exercise}
\begin{solution}
记 \(r=|x-a|\)、\(\sigma_1=2\pi\)、\(\sigma_2=4\pi\)。两种维数下径向导数均为
\[
 \Phi_{d,a}'(r)=-\frac1{\sigma_{d-1}r^{d-1}}.
\]
再求一次导数，便有 \(\Phi_{d,a}''+(d-1)\Phi_{d,a}'/r=0\)，所以原点移至 \(a\) 后，核仍在极点外调和。在球的外法向 \(n=e_r\) 下，通量为
\[
 \int_{\partial B(a,\varepsilon)}\partial_n\Phi_{d,a}\,\mathrm dS
 =-\frac1{\sigma_{d-1}\varepsilon^{d-1}}
      \cdot\sigma_{d-1}\varepsilon^{d-1}=-1.
\]
因此带负号的通量等于 \(1\)。对删球后的区域，穿过内边界向区域外走是朝向被删去的球心，故它的外法向为 \(-e_r\)，不是 \(e_r\)。

核在极点附近绝对可积：二维只需 \(\int_0^1r|\log r|\,\mathrm dr<\infty\)，三维只需 \(\int_0^1r\,\mathrm dr<\infty\)。乘上 \(\chi_R\) 后支集包含于 \(\overline{B(a,2R)}\)，故 \(K\in L^1\)。\(\chi_R\) 的所有正阶导数在 \(B(a,R)\) 内及 \(B(a,2R)\) 外都为零，因此 \(q\) 只在远离极点的壳层上可能非零，在那里各因子均光滑。补零后 \(q\) 全局光滑且紧支集。极点外使用普通乘积求导公式，得 \(-\Delta K=q\)；点质量则必须另从边界极限中算出。

取任意 \(\varphi\in C_c^\infty(\mathbb R^d;\mathbb C)\)，选足够大的外球包含 \(K\) 和 \(\varphi\) 的支集，再删去 \(\overline{B(a,\varepsilon)}\)，其中 \(0<\varepsilon<\min\{R,1\}\)。外球边界的项全为零。内球附近 \(K=\Phi_{d,a}\)，并且剩余区域的法向为 \(-e_r\)，Green 公式于是给出
\[
 \begin{aligned}
 &\int_{|x-a|>\varepsilon}K(x)(-\Delta\varphi(x))\,\mathrm dx\\
 &\quad=\int_{\mathbb R^d}q(x)\varphi(x)\,\mathrm dx
       +\int_{|x-a|=\varepsilon}
        \left(-\varphi\,\partial_r\Phi_{d,a}
                     +\Phi_{d,a}\partial_r\varphi\right)\,\mathrm dS.
 \end{aligned}
\]
第一项已可在全空间积分，因为 \(q\) 支撑于 \(|x-a|\geq R\)。边界上含 \(-\partial_r\Phi_{d,a}\) 的项正是 \(\varphi\) 的球面平均，趋于 \(\varphi(a)\)。另一项的模由 \(\|\nabla\varphi\|_\infty\) 乘 \(C\varepsilon|\log\varepsilon|\)（二维）或 \(C\varepsilon\)（三维）控制，故趋零。左侧由 \(K\) 的局部可积性趋于 \(\int K(-\Delta\varphi)\)，得到 \(-\Delta K=\delta_a+q\) 在紧支集测试函数上的等式。\(K\in L^1\)、\(q\in C_c^\infty\) 和 \(\delta_a\) 都定义温和分布，再由 \(C_c^\infty\) 在 \(\mathcal S\) 中稠密，将等式推广到 \(\mathcal S'\)。

最后取 \(\eta\in C_c^\infty(\mathbb R^d;\mathbb R)\)，在 \(\overline{B(a,2R)}\) 的一个邻域内等于 \(1\)。\(\Delta\eta\) 在 \(K\) 的支集上为零，而 \(\eta\) 在 \(q\) 的支集上及 \(a\) 点均为 \(1\)，故
\[
 0=\langle-\Delta K,\eta\rangle
  =\eta(a)+\int_{\mathbb R^d}q(x)\eta(x)\,\mathrm dx
  =1+\int_{\mathbb R^d}q.
\]
于是余项的总质量为 \(-1\)，恰好抵消点质量的总质量。空间截断保留了极点附近的单位通量，却必须在截断变化的壳层中产生这项补偿。
\end{solution}
